\section*{\fontsize{16pt}{15pt}\selectfont BUILDING A CAMPUS NETWORK \\USING SDN AND SD-WAN \\ABSTRACT}

\begin{itemize}
    \item[-] \textbf{Research issues:}

    Universities with several campuses need a network that is highly available, secure, centrally managed and able to connect all sites reliably. Traditional networks are configured device by device, lack a global view of the network and depend on a single WAN link. This project designs, implements and verifies a multi-campus network that combines a hierarchical campus design with Software-Defined Networking (SDN) and Software-Defined WAN (SD-WAN). The report includes 5 chapters:
        \begin{itemize}
        \item[•] Chapter 1 --- Introduction and Project Overview
        \item[•] Chapter 2 --- Theoretical Background
        \item[•] Chapter 3 --- Proposed Campus Network Design
        \item[•] Chapter 4 --- Implementation, Testing and Evaluation
        \item[•] Chapter 5 --- Conclusion and Future Work
        \end{itemize}
    \item[-] \textbf{Approach:}
        \begin{itemize}
        \item[•] Theory: Study VLAN, RSTP, EtherChannel, VRRP, OSPF, BGP, stateful firewalls and DMZ, OpenFlow 1.3, Open vSwitch, Ryu, ONOS and the Cisco SD-WAN (Viptela) architecture with OMP, TLOC, BFD, PKI and centralized policies.
        \item[•] Practice: Build a 70-node lab on EVE-NG with a main campus (Site 100), three branch campuses (Sites 200, 300, 400), an SD-WAN control site (Site 900), two simulated service providers (Internet and MPLS) and an expansion site (Site 500); verify every function layer by layer with automated measurements.
    \end{itemize}
    \item[-] \textbf{How to solve problems:}
    The main campus uses a Core--Distribution--Access design: two Cisco core switches provide VRRP gateways, OSPF and DHCP relay; an ASAv Active/Standby firewall pair protects the DMZ (Web, Mail) and the server farm (DHCP, DNS, Syslog). All six layer-2 switches are Open vSwitch instances controlled by a custom Ryu application (OpenFlow 1.3) that performs VLAN-aware MAC learning, a centrally computed forwarding tree with active link probing, a multi-path management channel (VLAN 99) independent of STP, self-healing after link or switch failures, VLAN and policy provisioning through an API, and a web-based Campus SDN Console. Branches use a Firewall-as-Core design. The sites are connected by Cisco SD-WAN 20.10.1 with eight vEdge routers over two transports, a BGP underlay, PKI authentication, IPsec overlay and vSmart policies (SLA classes, Application-Aware Routing, data policy and a control policy preferring each site's vEdge1); branch firewalls track the MPLS path with IP SLA to fail over to the Internet transport, and vManage sends e-mail alarms when an edge loses connectivity. Users at every site reach the Internet locally (Direct Internet Access) through NAT on the vEdge routers with an MPLS backup path, while the management zones and the controller site only have controlled Internet access.
    \item[-] \textbf{Results achieved:}
    Of 29 test cases, 25 passed and 3 partially. API-based VLAN and campus-wide policy changes took 3.8 and 8.0 ms, respectively. All six SDN link failures recovered within 3.5 seconds, and all ten vEdges joined with BFD sessions up. Ten counter-based tests confirmed SD-WAN policies. Branch traffic failed over from MPLS to the Internet in about 8--9 seconds (Site 200 re-measurement); hosts also accessed the Internet and resolved external names. Limitations include a single SDN controller, no high-load testing and no failover test at Site 500.
\end{itemize}

\newpage
